% GENERATED FILE. Edit canonical lessons, then run npm run generate:books.
% canonical-source-hash: fnv1a64:4923d2d9f680f92e
% canonical-lessons: TE-C218-phalitam, TE-C218-laksyam, TE-C218-udaharana, TE-C218-teda, TE-C218-srama

\chapter{Result, Goal, Example, Difference, Effort}
\label{ch:te-result-goal-example-difference-effort}
\hlchaptermodality{218}

\begin{chapteropening}
\textbf{By the end of this chapter:} \emph{I can say a result, a goal, an example, a difference and effort.}

Complete the last lesson of chapter 218: I can say a result, a goal, an example, a difference and effort.
\end{chapteropening}

\section[\texorpdfstring{\te{ఫలితం} (phalitaṁ)}{ఫలితం (phalitaṁ)}]{\te{ఫలితం} (phalitaṁ) — a result}

\label{lesson:TE-C218-phalitam}



\begin{quote}
Before the new one: say the Telugu for a custom, then the Telugu for a habit.
\end{quote}

\subsection*{You'll want to know: \te{ఫలితం}}
\textbf{\te{ఫలితం}} — \emph{phalitaṁ} — \textquotedblleft{}a result\textquotedblright{}.

\begin{grammarlens}[title={What you've built}]
One more word for this chapter.
\end{grammarlens}

\subsection*{Your turn}
\begin{itemize}
\raggedright
  \item \emph{Say it:} \emph{phalitaṁ}
  \item \emph{Say it:} \emph{phalitaṁ}, once more
  \item \emph{Say it:} say \emph{phalitaṁ}
  \item \emph{Recall:} say the Telugu for a custom, then the Telugu for a habit, then say \emph{phalitaṁ} again
\end{itemize}

\begin{tcolorbox}[breakable,colback=teal!4,colframe=teal!35!black,title={Before you move on}]
What does \te{ఫలితం} mean? (\textquotedblleft{}A result\textquotedblright{}.) Say it once more.
\end{tcolorbox}
\section[\texorpdfstring{\te{లక్ష్యం} (lakṣyaṁ)}{లక్ష్యం (lakṣyaṁ)}]{\te{లక్ష్యం} (lakṣyaṁ) — a goal, an aim}

\label{lesson:TE-C218-laksyam}



\begin{quote}
Before the new one: say the Telugu for a habit, then the Telugu for a result.
\end{quote}

\subsection*{You'll want to know: \te{లక్ష్యం}}
\textbf{\te{లక్ష్యం}} — \emph{lakṣyaṁ} — \textquotedblleft{}a goal, an aim\textquotedblright{}.

\begin{grammarlens}[title={What you've built}]
One more word for this chapter.
\end{grammarlens}

\subsection*{Your turn}
\begin{itemize}
\raggedright
  \item \emph{Say it:} \emph{lakṣyaṁ}
  \item \emph{Say it:} \emph{lakṣyaṁ}, once more
  \item \emph{Say it:} say \emph{lakṣyaṁ}
  \item \emph{Recall:} say the Telugu for a habit, then the Telugu for a result, then say \emph{lakṣyaṁ} again
\end{itemize}

\begin{tcolorbox}[breakable,colback=teal!4,colframe=teal!35!black,title={Before you move on}]
What does \te{లక్ష్యం} mean? (\textquotedblleft{}A goal, an aim\textquotedblright{}.) Say it once more.
\end{tcolorbox}
\section[\texorpdfstring{\te{ఉదాహరణ} (udāharaṇa)}{ఉదాహరణ (udāharaṇa)}]{\te{ఉదాహరణ} (udāharaṇa) — an example}

\label{lesson:TE-C218-udaharana}



\begin{quote}
Before the new one: say the Telugu for a result, then the Telugu for a goal.
\end{quote}

\subsection*{You'll want to know: \te{ఉదాహరణ}}
\textbf{\te{ఉదాహరణ}} — \emph{udāharaṇa} — \textquotedblleft{}an example\textquotedblright{}.

\begin{grammarlens}[title={What you've built}]
One more word for this chapter.
\end{grammarlens}

\subsection*{Your turn}
\begin{itemize}
\raggedright
  \item \emph{Say it:} \emph{udāharaṇa}
  \item \emph{Say it:} \emph{udāharaṇa}, once more
  \item \emph{Say it:} say \emph{udāharaṇa}
  \item \emph{Recall:} say the Telugu for a result, then the Telugu for a goal, then say \emph{udāharaṇa} again
\end{itemize}

\begin{tcolorbox}[breakable,colback=teal!4,colframe=teal!35!black,title={Before you move on}]
What does \te{ఉదాహరణ} mean? (\textquotedblleft{}An example\textquotedblright{}.) Say it once more.
\end{tcolorbox}
\section[\texorpdfstring{\te{తేడా} (tēḍā)}{తేడా (tēḍā)}]{\te{తేడా} (tēḍā) — a difference}

\label{lesson:TE-C218-teda}



\begin{quote}
Before the new one: say the Telugu for a goal, then the Telugu for an example.
\end{quote}

\subsection*{You'll want to know: \te{తేడా}}
\textbf{\te{తేడా}} — \emph{tēḍā} — \textquotedblleft{}a difference\textquotedblright{}.

\begin{grammarlens}[title={What you've built}]
One more word for this chapter.
\end{grammarlens}

\subsection*{Your turn}
\begin{itemize}
\raggedright
  \item \emph{Say it:} \emph{tēḍā}
  \item \emph{Say it:} \emph{tēḍā}, once more
  \item \emph{Say it:} say \emph{tēḍā}
  \item \emph{Recall:} say the Telugu for a goal, then the Telugu for an example, then say \emph{tēḍā} again
\end{itemize}

\begin{tcolorbox}[breakable,colback=teal!4,colframe=teal!35!black,title={Before you move on}]
What does \te{తేడా} mean? (\textquotedblleft{}A difference\textquotedblright{}.) Say it once more.
\end{tcolorbox}
\section[\texorpdfstring{\te{శ్రమ} (śrama)}{శ్రమ (śrama)}]{\te{శ్రమ} (śrama) — effort, toil}

\label{lesson:TE-C218-srama}



\begin{quote}
Before the new one: say the Telugu for an example, then the Telugu for a difference.
\end{quote}

\subsection*{You'll want to know: \te{శ్రమ}}
\textbf{\te{శ్రమ}} — \emph{śrama} — \textquotedblleft{}effort, toil\textquotedblright{}.

\begin{grammarlens}[title={What you've built}]
One more word for this chapter.
\end{grammarlens}

\subsection*{Your turn}
\begin{itemize}
\raggedright
  \item \emph{Say it:} \emph{śrama}
  \item \emph{Say it:} \emph{śrama}, once more
  \item \emph{Say it:} say \emph{śrama}
  \item \emph{Recall:} say the Telugu for an example, then the Telugu for a difference, then say \emph{śrama} again
\end{itemize}

\begin{tcolorbox}[breakable,colback=teal!4,colframe=teal!35!black,title={Before you move on}]
What does \te{శ్రమ} mean? (\textquotedblleft{}Effort, toil\textquotedblright{}.) Say it once more.
\end{tcolorbox}
